\section{Memory Analysis for BMS Subsystems}
This section examines the memory usage of BMS subsystems on the S32K358 microcontroller, considering both static random access memory (SRAM) and flash. Section XII-A summarizes the roles of BMS components on the S32K358 microcontroller, with SRAM handling runtime data and computations, while flash storing persistent code and calibration data. Finally, Section XII-B details memory allocation across specific subsystems, including SoC, SoP, CMU, and BJB.

\subsection{General Overview of BMS Components}
The BMS software on the S32K358 microcontroller utilizes both SRAM and flash, each serving distinct purposes. SRAM, a high-speed volatile memory, supports runtime computations by buffering cell voltages, temperatures, pack current, and status flags \cite{memoryDefinitions}. Flash, on the other hand, is non-volatile, storing persistent data such as executable code, initialization routines, drivers, lookup tables, configuration parameters and calibration data, ensuring critical information is retained even when power is removed \cite{memoryDefinitions}. Figure~10 illustrates the memory usage of the BMS algorithms, which occupy 31.83\% of SRAM and 5.29\% of flash, with SRAM primarily consumed by computational tasks and flash dominated by drivers and initialization routines rather than algorithm code.
In this work, the largest allocation to SRAM, 57.05\%, is assigned to SoC estimation because 105 cells are continuously tracked via Kalman filters. The CMU follows at 18.59\%, buffering a total of 245 voltage/temperature channels, while system-level logic consumes 13.28\% for core control and task scheduling. The SoP estimation occupies 9.79\%, providing continuous power and operational current limits. Lastly, lightweight components such as passive cell balancing, BJB, and fault management each use less than \textasciitilde{1\%} of SRAM: the cell balancing only toggles bleed resistors, the BJB handles aggregate pack-level measurements rather than detailed per-cell data, and fault management applies simple threshold checks.

On the other hand, flash is mostly consumed by system‑level functions at 72.08\%, for drivers, communication protocols, and initialization routines. Next, the SoP consumes 10.18\%, reflecting storage for its lookup tables, thresholds, and calibration data needed to compute charge/discharge current and power limits. In comparison, the SoC consumes 7.77\%, mainly storing estimation algorithms and per-cell calibration constants required for accurate state-of-charge computation across 105 cells. Furthermore, CMU occupies 6.98\% because it primarily performs real-time measurements and preprocessing, storing minimal persistent data. Passive cell balancing, BJB and fault management each consume less than 2\% of the used flash due to their focused functionality—balancing logic, limited high-voltage measurements (a total of 5 HV measurements), and threshold-based fault detection—each requiring only lightweight processing and data storage.
\subsection{Specific Analysis Of Algorithms and Modules}
To further complement the previous analysis, a detailed breakdown of SoC and SoP algorithms, and CMU and BJB modules is provided in this section. In the SoC algorithm, each of the 7 temperature-specific EKFs consumes an equal 14.29\% share of its total SRAM usage, collectively accounting for 57.05\% of the algorithm’s footprint, due to their identical structure and parallel execution. In flash, the distribution is similarly uniform, with each EKF using approximately 14.30\%, reflecting consistent model storage across temperature-specific estimators.

For SoP, SRAM usage is led by power limit computations at 50.45\%, which involve storing intermediate results and thresholds derived from lookup tables; the remainder percentage is used for current limit functions, requiring less memory due to their simpler structure. Conversely, in flash memory, current limit functions occupy 59.54\%, due to its persistent code requirements and 16 2-D lookup tables, half for charge and half for discharge, while power limit routines use comparatively less non-volatile space with 10 2-D lookup tables in total, evenly split between charge and discharge.
The CMU module, built around the MC33775A IC, allocates 88.73\% of its total SRAM usage to temperature conversions across 140 channels due to the complexity of NTC math and fault logic, while the remaining memory supports faster but simpler cell voltage computations. In flash, cell and pack voltage computations take the larger share at 41.75\%, due to their broader data handling and communication logic across 105 channels, while temperature conversions—though applied to more channels—require less storage owing to their uniform processing and slower dynamics.
Finally, the BJB module, built around the MC33772C IC, EDM and DCFC voltage computations consume 57.88\% of its total SRAM usage due to diagnostic, stability monitoring and safety-critical routines; in comparison, precharge resistor temperature sensing accounts for 21.23\%, as continuous monitoring is required for real-time protection against thermal risks during precharge. The remainder is used for pack current sensing, which requires fewer resources due to its lower runtime complexity. In flash, EDM and DCFC voltage sensing dominate at 54.51\% for persistent diagnostics, compliance and long-term measurement accuracy; whereas precharge resistor temperature uses 28.82\% for thermal modelling. By contrast, the code for pack current takes the smallest portion, justified by its simpler algorithmic structure.
\newpage

\begin{figure*}[htbp]
  \centering
  \includegraphics[width=\textwidth]{Figures/Fig78.png}
  \caption{SRAM and flash memory consumption by different BMS algorithms and modules.}
  \label{Fig78}
\end{figure*}
